\section{会话管理与上下文控制}
%% ============================================================

Claude Code 的输出质量与上下文状态密切相关。管理好会话是长时间使用的关键。

\subsection{不同任务之间用 /clear 清空（Tips \#12, \#24）}

\begin{importantbox}{最重要的会话管理原则}
一个干净会话配合精准 prompt，胜过一个混乱的三小时长会话。不同任务？先 \texttt{/clear}。当你在同一个问题上纠正 Claude 超过两次，也应该 \texttt{/clear} 重新开始——上下文中堆积的失败尝试会"污染"后续输出。
\end{importantbox}

\subsection{扩展上下文窗口到 1M Tokens（Tip \#10）}

Sonnet 4.6 和 Opus 4.6 都支持 1M token 上下文窗口。在 Max/Team/Enterprise 计划中，Opus 自动升级到 1M。也可以在会话中用 \texttt{/model opus[1m]} 切换。

相关环境变量：
\begin{itemize}
  \item \texttt{CLAUDE\_CODE\_AUTO\_COMPACT\_WINDOW}：控制自动 compaction 触发时机
  \item \texttt{CLAUDE\_AUTOCOMPACT\_PCT\_OVERRIDE}：设置 compaction 百分比阈值
\end{itemize}

\subsection{引导 Compaction 保留关键信息（Tip \#21）}

当上下文自动压缩时，用命令引导 Claude 保留什么：

\begin{lstlisting}[caption={引导 compaction 示例}, language=bash]
/compact focus on the API changes and the list of modified files
\end{lstlisting}

也可以在 CLAUDE.md 中添加常驻指令："When compacting, preserve the full list of modified files and current test status."

\subsection{用 /btw 快速提问（Tip \#14）}

\texttt{/btw} 弹出一个覆盖层用于快速提问，不会进入对话历史。适合问"为什么选了这个方案？"或"另一个选项的权衡是什么？"主上下文保持干净。

\subsection{用 /branch 尝试不同方案（Tip \#43）}

\texttt{/branch}（或 \texttt{/fork}）在当前节点创建对话副本。在分支中尝试激进重构，如果成功就保留，失败了原始对话不受影响。与 \texttt{Esc+Esc} 回滚不同，两条路径同时存活。

\subsection{用 /loop 做周期性检查（Tip \#22）}

\begin{lstlisting}[caption={/loop 示例}, language=bash]
/loop 5m check if the deploy succeeded and report back
/loop 20m /review-pr 1234
\end{lstlisting}

支持 \texttt{s}, \texttt{m}, \texttt{h}, \texttt{d} 时间单位，默认 10 分钟。任务绑定到会话，3 天后过期，不会无限运行。适合监控部署、CI 流水线等场景。

\subsection{本章小结}

会话管理的关键：不同任务间 \texttt{/clear}、修正两次仍失败就重新开始、引导 compaction 保留关键信息、用 \texttt{/btw} 和 \texttt{/branch} 保持上下文干净。理解上下文窗口和 compaction 机制是高效使用 Claude Code 的基础。

%% ============================================================
